\section{Results}

In the following results, item means, option frequencies, and upper score range proportions are derived from the questionnaire platform's aggregate export table; internal consistency coefficients, exploratory correlations, and group comparisons are based on the subsequently compiled participant-level response matrix ($N = 32$).

\subsection{Sample Characteristics}

The core questionnaire sample consisted of 32 female college students.
The majority were aged 18--20 (28 participants, 87.50\%).
Regarding exercise frequency in the past month, 15 participants (46.88\%) exercised 2--3 times per week, 6 (18.75\%) exercised once per week, 6 (18.75\%) exercised 4 or more times per week, and 5 (15.63\%) rarely exercised.
Twenty-five participants (78.13\%) had used the campus gym, while 7 (21.88\%) had not.

Gym area usage showed marked differences.
The most frequently selected area was the cardio area with treadmills, ellipticals, and similar equipment (21 participants, 65.63\%), followed by the fixed machine area (9 participants, 28.13\%).
Only 5 participants (15.63\%) selected the free-weight area, and 8 participants (25.00\%) reported that they generally did not enter the gym.
Strength training experience was also limited: 10 participants (31.25\%) had no exposure at all, 14 (43.75\%) had tried it but not systematically, and only 5 (15.63\%) were currently engaged in regular strength training.
Regarding exercise goals, fat loss/slimming was the most common (18 participants, 56.25\%), followed by toning/sculpting and improving fitness at 11 participants each (34.38\%), and building strength at 6 participants (18.75\%).

\begin{table}[H]
\centering
\caption{Sample Exercise Experience and Gym Usage Overview ($N = 32$)}
\label{tab:sample_actual}
\begin{tabular}{lll}
\toprule
Variable & Main Result & Percentage \\
\midrule
Has used campus gym & 25 participants & 78.13\% \\
Exercises 2--3 times per week & 15 participants & 46.88\% \\
Most frequently uses cardio area & 21 participants & 65.63\% \\
Most frequently uses free-weight area & 5 participants & 15.63\% \\
No strength training experience at all & 10 participants & 31.25\% \\
Some unsystematic attempts & 14 participants & 43.75\% \\
Currently engaged in regular strength training & 5 participants & 15.63\% \\
Primary goal is fat loss/slimming & 18 participants & 56.25\% \\
Primary goal is building strength & 6 participants & 18.75\% \\
\bottomrule
\end{tabular}
\end{table}

Figure \ref{fig:gym_area_use} more intuitively shows that the free-weight area accounts for a markedly lower proportion of commonly used areas compared to the cardio area.

\begin{figure}[H]
\centering
\includegraphics[width=0.85\textwidth]{gym_area_use.png}
\caption{Most Frequently Used Gym Areas}
\label{fig:gym_area_use}
\end{figure}

\subsection{Spatial Pressure and Free-Weight Avoidance}

The spatial pressure/training anxiety construct comprised Items 10--15, with an individual-level mean of 3.72 ($SD = 0.77$) and good internal consistency (Cronbach's $\alpha = .81$).
The most prominent item was ``I feel noticeably uncomfortable when there are many men in the free-weight area'' ($M = 4.24$, upper score range 81.25\%), followed by ``The population density, equipment concentration, and crowding in the free-weight area make me not want to stay'' ($M = 3.88$, upper score range 75.00\%) and ``I worry that if my form is incorrect, I will be noticed or judged by others'' ($M = 3.81$, upper score range 65.63\%).
This indicates that participants' discomfort does not stem solely from abstract body anxiety, but is concentrated in specific spatial conditions such as male predominance, open visibility, form errors, and crowded conditions that discourage lingering.

\begin{table}[H]
\centering
\small
\caption{Spatial Pressure/Training Anxiety Item Summary}
\label{tab:spatial_items_actual}
\begin{tabularx}{\textwidth}{lXcc}
\toprule
Item & Brief Description & Mean & Upper Score Range \\
\midrule
Q10 & Overall layout of the free-weight area feels inaccessible & 3.19 & 40.63\% \\
Q11 & Noticeably uncomfortable when many men are present & 4.24 & 81.25\% \\
Q12 & Mirrors and openly visible environment heighten sense of being evaluated & 3.57 & 53.13\% \\
Q13 & Worry about being noticed or judged for incorrect form & 3.81 & 65.63\% \\
Q14 & Free-weight area feels more like a male-dominated space & 3.61 & 62.50\% \\
Q15 & Population density, equipment concentration, and crowding reduce willingness to stay & 3.88 & 75.00\% \\
\bottomrule
\end{tabularx}
\end{table}

The participant-level data further support this interpretation.
After coding responses from ``never'' to ``always'' as 1--5, the mean frequency of actively avoiding the free-weight area was 3.09 ($SD = 1.33$).
Spatial pressure/training anxiety was moderately to strongly positively correlated with free-weight area avoidance frequency (Spearman's $\rho = .65$, $p < .001$).
Spatial pressure was also positively correlated with the intervention preference index ($\rho = .67$, $p < .001$), indicating that participants who felt greater spatial pressure in the free-weight area were more inclined to support measures such as workshops, semi-private partitioning, and stereotype reduction.

\begin{table}[H]
\centering
\small
\caption{Individual-Level Construct Descriptives and Reliability ($N = 32$)}
\label{tab:participant_scales}
\begin{tabular}{lcccc}
\toprule
Construct/Item Set & No. of Items & Mean & $SD$ & Cronbach's $\alpha$ \\
\midrule
Spatial pressure/training anxiety & 6 & 3.72 & 0.77 & .81 \\
Social media appearance internalization & 3 & 2.66 & 0.93 & .76 \\
Training self-efficacy & 1 & 3.49 & 0.94 & -- \\
Intervention preference index & 3 & 3.98 & 0.54 & N/A \\
\bottomrule
\end{tabular}
\end{table}

\begin{table}[H]
\centering
\small
\caption{Key Exploratory Spearman Correlations}
\label{tab:spearman_results}
\begin{tabular}{lcc}
\toprule
Variable Pair & $\rho$ & $p$ \\
\midrule
Spatial pressure/training anxiety -- Free-weight area avoidance frequency & .65 & $< .001$ \\
Spatial pressure/training anxiety -- Intervention preference index & .67 & $< .001$ \\
Intervention preference index -- Free-weight area avoidance frequency & .45 & .009 \\
Social media appearance internalization -- Free-weight area avoidance frequency & .31 & .084 \\
Training self-efficacy -- Free-weight area avoidance frequency & -.11 & .562 \\
\bottomrule
\end{tabular}
\end{table}

Figure \ref{fig:spatial_anxiety} shows the score distribution for each item of the spatial pressure/training anxiety construct.
The upper score range was most concentrated for items regarding male predominance, crowded conditions, and fear of incorrect form.

\begin{figure}[H]
\centering
\includegraphics[width=\textwidth]{slider_spatial_anxiety.png}
\caption{Spatial Pressure/Training Anxiety Slider Distribution}
\label{fig:spatial_anxiety}
\end{figure}

\begin{figure}[H]
\centering
\includegraphics[width=0.78\textwidth]{survey_spearman_correlations.png}
\caption{Exploratory Spearman Correlation Matrix}
\label{fig:spearman_matrix}
\end{figure}

The frequency of actively avoiding the free-weight area also supports this picture.
Five participants (15.63\%) reported never avoiding, 5 (15.63\%) rarely avoiding, 10 (31.25\%) sometimes avoiding, and 6 (18.75\%) often avoiding.
In the aggregate export table, the ``always'' option was blank; however, since the valid sample for this item was 32 and the other options summed to 26 participants, ``always'' could be inferred as 6 participants (18.75\%).
In other words, at least 22 participants (68.75\%) sometimes or more frequently actively avoided the free-weight area.

Group comparisons also revealed differences between participants who actually entered the free-weight area and those who did not.
Only 5 participants most frequently used the free-weight area, so these results should be interpreted only exploratorily; however, this group's mean spatial pressure/training anxiety score was lower than that of non-regular users (2.60 vs.\ 3.92, Mann--Whitney $U = 16.00$, $p = .008$), and their avoidance frequency was also lower (1.40 vs.\ 3.41, $U = 10.00$, $p = .002$).
Participants with systematic strength training experience (currently or previously engaging in regular strength training, $n = 8$) also had lower avoidance frequency than those without systematic experience (1.88 vs.\ 3.50, $U = 28.00$, $p = .003$).
These differences cannot demonstrate that training experience directly reduces avoidance, but they suggest that familiarity with equipment and training procedures may co-occur with lower spatial pressure and avoidance tendencies.

Among the reasons for avoidance, the most frequently reported were ``too many men'' and ``do not know how to use equipment,'' both selected by 19 participants (59.38\%); followed by ``do not know where to start'' (16 participants, 50.00\%), ``no companion'' (11 participants, 34.38\%), ``fear of making form errors'' (9 participants, 28.13\%), and ``layout/atmosphere is too intimidating'' (8 participants, 25.00\%).
Notably, ``do not want to build visible muscle'' was selected by 0 participants (0.00\%).
Together with the social media items discussed below, this indicates that avoidance in this sample was not primarily driven by fear of becoming too muscular, but rather by a combination of spatial belonging, equipment knowledge, and evaluation pressure as a beginner.

\begin{figure}[H]
\centering
\includegraphics[width=0.85\textwidth]{avoidance_reasons.png}
\caption{Primary Reasons for Avoiding the Free-Weight Area (only options with explicit frequencies in the aggregate export table are shown)}
\label{fig:avoidance_reasons}
\end{figure}

\subsection{Social Media and Body Ideals}

Social media use was highly prevalent in the sample.
Xiaohongshu (RED) was the most frequently used platform (30 participants, 93.75\%), followed by Bilibili (22 participants, 68.75\%), Weibo (9 participants, 28.13\%), and Douyin/TikTok (8 participants, 25.00\%).
However, participants generally did not spend much time daily browsing body, fitness, fat loss, appearance, or outfit-related content: 14 participants (43.75\%) spent less than 15 minutes, 16 (50.00\%) spent 15--30 minutes, 2 (6.25\%) spent 31--60 minutes, and none selected more than 1 hour.

The social media appearance internalization construct comprised Items 21--23, with an individual-level mean of 2.66 ($SD = 0.93$) and acceptable internal consistency (Cronbach's $\alpha = .76$).
Among these items, ``Appearance-related content on social media influences what kind of exercise I choose'' had a mean of 3.40 and an upper score range of 59.38\%, indicating that appearance content does enter the exercise selection process.
However, ``I worry that strength training will make me look not slender enough or not conform to mainstream feminine aesthetics'' had a mean of only 1.85, with 81.25\% in the lower score range and only 15.63\% in the upper score range.
``I treat the female body standards popular on social media as my own goal'' had a mean of 2.73 and an upper score range of 31.25\%.
Thus, media influence in this study functioned more as a limited background factor in exercise selection rather than as a strong, unidirectional feminine aesthetic pressure.
Correlation analysis also supports this cautious interpretation: the correlation between social media appearance internalization and free-weight area avoidance frequency was positive but did not reach conventional significance ($\rho = .31$, $p = .084$), which was weaker than the relationship between spatial pressure/training anxiety and avoidance frequency.

\begin{table}[H]
\centering
\small
\caption{Social Media Appearance-Related Item Summary}
\label{tab:media_items_actual}
\begin{tabularx}{\textwidth}{lXccc}
\toprule
Item & Brief Description & Mean & Lower Score Range & Upper Score Range \\
\midrule
Q21 & Appearance content influences exercise choice & 3.40 & 21.88\% & 59.38\% \\
Q22 & Worry that strength training will not look slender or conform to feminine aesthetics & 1.85 & 81.25\% & 15.63\% \\
Q23 & Internalize popular social media female body standards & 2.73 & 43.75\% & 31.25\% \\
\bottomrule
\end{tabularx}
\end{table}

\begin{figure}[H]
\centering
\includegraphics[width=0.85\textwidth]{slider_media_internalization.png}
\caption{Social Media Appearance-Related Slider Distribution}
\label{fig:media_distribution}
\end{figure}

\subsection{Self-Efficacy and Intervention Preferences}

Item 24 showed that participants' self-efficacy for achieving their training goals was at a moderate-to-upper level ($M = 3.49$, upper score range 50.00\%).
This suggests that many participants did not completely dismiss their ability to participate in gym training, but this confidence was insufficient to overcome the spatial and beginner barriers associated with the free-weight area.
Corroborating this interpretation, training self-efficacy showed virtually no association with free-weight area avoidance frequency ($\rho = -.11$, $p = .562$), indicating that general confidence in achieving training goals did not directly translate into willingness to enter the free-weight area.

The intervention preference index had the highest mean among all item sets, with an individual-level mean of 3.98 ($SD = 0.54$).
Because the three items measured three distinct interventions---media stereotype reduction, semi-private space, and beginner-friendly women's strength training workshops---this index was not interpreted as a strict psychological scale but rather as a descriptive indicator of overall support.
Among these items, ``If the university offered a beginner-friendly women's strength training workshop, I would be more willing to try strength training'' had a mean of 4.50 and an upper score range of 96.88\%.
``If semi-private partitioning or visual screening were added to the free-weight area, I would be more willing to use it'' had a mean of 4.04 and an upper score range of 81.25\%.
``If stereotypes about women's strength training on social media decreased, I would be more willing to participate in strength training'' had a mean of 3.39 and an upper score range of 50.00\%.
This ordering indicates that the most desired support was first university-based introductory guidance and spatial adjustments, followed by broader media stereotype improvement.
Notably, the intervention preference index was moderately positively correlated with free-weight area avoidance frequency ($\rho = .45$, $p = .009$), suggesting that participants who avoided the free-weight area more frequently were also more inclined to support these intervention measures---the strongest demand for change came from those most affected by the barriers.

\begin{figure}[H]
\centering
\includegraphics[width=0.85\textwidth]{slider_intervention_acceptance.png}
\caption{Intervention Preference Slider Distribution}
\label{fig:intervention_distribution}
\end{figure}

\subsection{Qualitative Themes}

The interview material yielded five themes directly related to the questionnaire results.
First, the strength training area was experienced as a space of heightened visibility.
P1 explained the difference between the cardio area and the strength area in terms of bodily orientation and mirrors: during cardio, one has one's ``back to others,'' while the strength area ``has mirrors,'' making it easier to see others and to become aware of being seen oneself.
P3 also mentioned that during strength training, ``there is a large mirror right in front, and it is easy to notice surroundings in peripheral vision,'' occasionally feeling that someone was looking at her.

Second, beginner barriers and fear of incorrect form were central obstacles.
P3 said that part of the reason she avoided the strength area was that she ``did not know how to do anything,'' worrying about not knowing how to use the equipment and having non-standard form, which might lead others to think she was ``really bad at this.''
P2 offered a similar explanation from a more measured perspective: many female students ``do not really understand strength training methods'' or how to use the tools, so when it is crowded, they worry about having non-standard form and delaying others' use of equipment.
This is consistent with the high frequency of selecting ``do not know how to use equipment'' and ``do not know where to start'' in the questionnaire.

Third, the male-dominated atmosphere created a sense of distance, but the interviewees did not interpret it simply as hostility from individual men.
P2 considered the gym to be well-equipped and that people were generally respectful, but also noted that the space felt crowded when busy, and the gym ``mostly looks like it was designed in a masculine style,'' with most lifters being male; consequently, female students might feel their experience and weight levels differed from others, creating embarrassment.
P3 similarly emphasized that the male lifters were well-trained with considerable muscle mass---this was not itself a target of criticism, but it could create a sense of intimidation for beginners.

Fourth, avoidance often took the form of strategic negotiation rather than complete withdrawal from exercise.
P1 would switch to cardio, find a roommate to accompany her, or go during less crowded times before 9 a.m.
P3 would choose not to wear glasses and go with a friend, relying on ``blocking out the surroundings'' to calm herself down.
In the questionnaire, the most common coping strategies were also going during off-peak timing (46.88\%), staying only in the cardio/fixed machine area (37.50\%), and giving up on strength training they had originally planned or leaving the gym (both 31.25\%).

Fifth, social media appearance influence existed, but the three interviewees each had different relationships with media content, and none translated appearance pressure into direct avoidance of strength training.
P1 frequently encountered fitness videos but reported ``not really seeing'' negative appearance-related evaluations; when she saw fit bodies, she thought ``that is great'' and ``did not really care what people said in the comments''---her engagement with social media fitness content was more appreciative than anxious.
P2 simultaneously encountered both positive and negative viewpoints but judged that ``the positive views are a bit more'' and believed that ``mainstream aesthetics are quite diverse nowadays''; regarding the claim that ``getting too muscular does not look good,'' although she ``felt uncomfortable,'' she resolved the conflict by redefining strength training---``fitness is not necessarily about getting bulky, but about making your muscles and posture better'' and ``as long as body fat percentage goes down, it is fine.''
P3 most directly rejected appearance pressure, ``completely disagreeing'' that strength training would deviate from mainstream aesthetics and ``sometimes even arguing back''; however, she mentioned that social media content about ``bad people and incidents'' at gyms could influence her overall impression of the gym, suggesting that social media influence may not operate through the appearance anxiety pathway but rather indirectly reduce willingness to enter by disseminating negative gym culture.
These three differentiated response patterns corroborate the asymmetric pattern found in the questionnaire---Q21 (appearance content influences exercise choice) had an upper score range of 59.38\%, but Q22 (fear of not looking slender enough from strength training) had a lower score range of 81.25\%: appearance content does enter participants' informational environment, but its role is more akin to a limited background influence rather than a mechanism directly driving avoidance.

\begin{table}[H]
\centering
\small
\caption{Integration of Interview Themes and Quantitative Results}
\label{tab:qual_themes_actual}
\begin{tabularx}{\textwidth}{lX}
\toprule
Theme & Relationship to Questionnaire Results \\
\midrule
Spatial visibility and mirror gaze & Explains Q12 upper score range of 53.13\% and interview accounts that the strength area produces a greater sense of being seen compared to the cardio area \\
Beginner barriers and fear of incorrect form & Corresponds to ``do not know how to use equipment'' 59.38\%, ``do not know where to start'' 50.00\%, and Q13 upper score range of 65.63\% \\
Male-dominated atmosphere and sense of distance & Corresponds to ``too many men'' 59.38\%, Q11 upper score range of 81.25\%, and Q14 upper score range of 62.50\% \\
Strategic negotiation rather than complete withdrawal & Corresponds to off-peak timing, switching to cardio/fixed machine area, and giving up planned strength training as coping strategies \\
Limited and diverse appearance influence & Corresponds to Q21 showing influence but Q22 scoring low, indicating media influence exists but is not a primary avoidance mechanism \\
Beginner-friendly/women-friendly intervention needs & Corresponds to beginner-friendly women's strength training workshop upper score range of 96.88\% and semi-private partitioning upper score range of 81.25\% \\
\bottomrule
\end{tabularx}
\end{table}
